\chapter{Background and Related Work}\label{background-and-related-work}

This section builds towards an architecture that can answer a policy questionnaire from the open web without asserting more than its evidence supports. It begins with the ODMI itself, followed by the criteria that must be met for it to remain a legitimate assessment post-automation. Past works are then weighed against those criteria. The review turns first to LLMs, why they are a sensible starting point and how they become agentic, then to the limitations that make them unreliable where evidence is scarce. Self-verification, multi-agent debate and abstention answer those limitations. The section closes with the difficulty a multi-lingual setting adds, and with the gap that remains for this work to address.

\section{The ODMI as an Assessment}\label{the-odmi-as-an-assessment}

\includegraphics[width=7.08611in,height=1.73819in]{/Users/benjyb/Desktop/MscProject/.claude/worktrees/odmi-figures-reconciliation-244489/Dissertation/latex/media/media/image2.png}

\ccnote{The four dimensions of the ODMI aim to measure} the entire open-data lifecycle, from creating policy to delivering real-world benefits. Policy measures open-data strategies, government structures, laws and implementation measures. Portal measures the functionality and usability of the national open data portal. Quality measures metadata standards, compliance, monitoring and data deployment quality. \ccnote{Impact measures the reuse of open data and the benefits it produces.} These dimensions perform sequential roles, with Policy informing decisions on open data, the Portal providing access to data, Quality assuring it is of a usable standard, and Impact measuring the value of this data.

A critique of the ODMI report is that it uses self-reporting by national open-data teams, which may lead to misinterpreted questions, reporting inconsistencies and self-reporting bias. The template for answers requests evidence but 15.1\% of the 5,148 pairs carry no comment at all, and if there is a populated comment it is frequently not evidence, such as a bare portal URL or one-word entries like ``yes''. Around 20\% of the questions cover internal practices (PT28 asks whether they monitor search key words, for example) and span all four dimensions. While Capgemini acts as an external verifier, they only change 10\% of the answers with counterevidence, complement the submitted answer with further evidence 27\% of the time, and confirm the original answer 63\% of the time.

Davis, Kingsbury and Merry (2012) argue that indicators, such as the ODMI, function as a technology of governance, ranking states, shaping standard-setting and influencing how resources are allocated and decisions are made. Their central concern is what reliance on indicators does to the authority and contestability of the decisions that follow. An indicator that has these influences must be open to scrutiny, and a country that disputes its ranking must be able to see the evidence upon which it was assessed. Contestability is the condition on which an indicator's legitimacy rests and is what a clear evidence trail supplies. On this reading, the ODMI's reliance on self-reporting means answers may rest on unevidenced or \ccnote{unverifiable assertions. 

Fumega and Gao (2026) put this directly, asking whether AI can take over indicator development, evidence discovery and policy evaluation without compromising `contextual nuance or democratic legitimacy' (p.1). Mirroring the results of human assessors is only half of the question. The other half is whether artificial intelligence can preserve the epistemic properties that make these assessments worth having, such as explainability.} Handing the job to a machine may cut costs, but if the underlying processes are opaque and indisputable, this technology ultimately degrades the authority and contestability of global benchmarking data \claudenote{(Davis et al., 2012)}.

\section{Criteria for an Automated System}\label{criteria-for-an-automated-system}

The criteria below set out what such a system has to demonstrate before it can be trusted with the automation of producing benchmarking indices. They govern any generative system that stands in for human judgement in producing a published assessment comparing several parties, where the system finds its own evidence. That covers the benchmarking indices in §2.1 and comparable exercises such as sustainability ratings and regulatory compliance monitoring.

The United Nations Economic Commission for Europe (UNECE) note that little specific guidance or established practice yet exists for using AI in generating official statistics and observe that most research remains focused on predictive rather than generative models (UNECE, 2025). No framework was identified during this work that lays out the criteria necessary for a reliable generative automation system. The criteria below are therefore assembled here, drawing on two established bodies of work with some adaptation to fit the setting.

Statistical agencies have for some time been automating work that people used to do by hand, and they have a settled way of deciding when that substitution is warranted. The Quality Framework for Statistical Algorithms (QF4SA) was published by the UNECE in 2022 (Yung et al., 2022) and was built with machine learning algorithms in mind. It directly addresses how previous statistical quality frameworks were written to judge a finished statistic, but not the algorithm that produced it, and they argue how an answer was reached is as important as the answer itself.

The QF4SA proposes five dimensions against which a statistical organisation should judge an algorithm before adopting it. Two of them, Cost and Timeliness, describe what an assessment costs to adopt and how quickly it runs, which are irrelevant to whether the process maintains legitimacy post-automation, so only three are carried forward here. Explainability is the degree to which a person can understand how an output follows from the inputs to it. Accuracy is the closeness of a computation to the value it was intended to measure. Reproducibility is the ability to obtain the same result from the same data and method. These dimensions are carried through, with explainability and accuracy extended into Attributability and Convergent Validity respectively, to reflect what they measure once the assessor is generative.

Williamson, Xi and Breyer (2012) set out the evidence an automated scorer must produce before it may be used operationally. Two of the criteria below come from their framework, Generalisability and Subgroup Equity, and a third shapes how accuracy is treated. \ccnote{Their framework was drafted for deterministic automated scoring of written answers.} It also assumes a human marks alongside the machine, and states that a machine marking alone would need stricter criteria than it currently supplies.

Convergent Validity asks whether the system reaches the right answer. The QF4SA defines accuracy as closeness to a true value, and \ccnote{§3.4} discusses the limits of self-reported answers as a stand-in for that value. Williamson et al. (2012) add that agreement with an independent measure carries more weight than agreement with a human key alone. Nine questions in the ODMI\textquotesingle s Quality dimension are deterministic, asking for a proportion of datasets meeting a metadata standard, and \ccnote{§3.5} sets out how these are computed directly.

\textbf{\ccnote{Table 2.1}: Criteria for a generative system automating an assessment.}

{\def\LTcaptype{none} % do not increment counter
\begin{longtable}[]{@{}
  >{\raggedright\arraybackslash}p{(\linewidth - 8\tabcolsep) * \real{0.1565}}
  >{\raggedright\arraybackslash}p{(\linewidth - 8\tabcolsep) * \real{0.1569}}
  >{\raggedright\arraybackslash}p{(\linewidth - 8\tabcolsep) * \real{0.2579}}
  >{\raggedright\arraybackslash}p{(\linewidth - 8\tabcolsep) * \real{0.2299}}
  >{\raggedright\arraybackslash}p{(\linewidth - 8\tabcolsep) * \real{0.1987}}@{}}
\toprule\noalign{}
\begin{minipage}[b]{\linewidth}\centering
\textbf{Criterion}
\end{minipage} & \begin{minipage}[b]{\linewidth}\centering
\textbf{What it asks}
\end{minipage} & \begin{minipage}[b]{\linewidth}\centering
\textbf{Why it matters here}
\end{minipage} & \begin{minipage}[b]{\linewidth}\centering
\textbf{What the system must show}
\end{minipage} & \begin{minipage}[b]{\linewidth}\centering
\textbf{Source, and what was adapted}
\end{minipage} \\
\midrule\noalign{}
\endhead
\bottomrule\noalign{}
\endlastfoot
Convergent Validity & Does the system reach the right answer? & The ODMI key is self-reported and can be a cycle old, so no single true value exists to score against (\ccnote{§3.4}). & Agreement with the ODMI key, and agreement with a measure computed independently of it. Nine Quality questions computed straight from catalogue metadata supply that second measure (\ccnote{§3.5}). & QF4SA accuracy (Yung et al., 2022). Williamson et al. (2012) add that an independent measure carries more weight than a human key alone. \\
Generalisability & Does performance hold across every kind of question the assessment requires? & An aggregate figure can be carried by the easy questions and hide failure on the rest. \claudenote{Running end to end over the open web, Fumega and Gao (2026) report mismatch against the human answers of 19\% on an indicator anchored in a single document against 69\% on one whose evidence is scattered, averaging 42\% over four indicators.} & Results decomposed by answer type and by ODMI dimension, not one headline figure. & Williamson et al. (2012). \\
Subgroup Equity & Are the countries being compared held to the same standard? & A country's score should follow from its open-data maturity, not from how searchable its web presence is or how well the model reads its language. & Adequate performance within each subgroup, not only overall, with results split by native language and maturity. & Williamson et al. (2012). Schenk and Kern (2024) extend QF4SA with algorithmic fairness, amending accuracy to require performance across subgroups. \\
Attributability & Does every committed answer trace to a passage a reader would agree supports it? & A generative model can produce a fluent answer with no basis (§2.4), and a country can only contest a score it is shown the evidence for. & Two conditions. The quoted passage exists in the cited source, and a reasonable reader agrees it supports the claim. & The QF4SA explainability (Yung et al., 2022), sharpened by the reader test in Rashkin et al. (2023). \\
Selectivity & Can the system decline to answer, and was it right to decline? & Models are optimised to guess when uncertain (Kalai et al., 2025). The QF4SA assumes the data is already available, but that is not the case for the ODMI, where the evidence has to be found and may not exist in public. & A commitment only where the evidence supports it, an abstention otherwise, and a stated confidence that tracks observed accuracy. & New here. Neither framework supplies it, as both assume a system that always answers. \\
Reproducibility & Does a second run of the assessment return the same answers? & A deterministic algorithm satisfies this by construction and a generative assessor cannot. Two runs may search differently while the truth has not changed, so the outcome is what has to hold, not the route. & Traceability, a property of the record: prompts, evidence, model version and settings, enough to re-derive any reported figure without re-running the system. Stability, a property of the system, measured as variance between runs. & QF4SA reproducibility (Yung et al., 2022), the one dimension carried through unchanged. \\
\end{longtable}
}

Together these six criteria\claudenote{, set out in Table 2.1,} make up the conditions under which an automated ODMI score remains contestable in Davis et al.\textquotesingle s sense. A country disputing its score should be able to check the evidence held against it (Attributability), and know the answers weren't guessed (Selectivity). It should also be able to establish whether the same standard was applied to its neighbours (Subgroup Equity), whether the result agrees with what a human assessor would have reached (Convergent Validity), whether the score rests on the parts of the questionnaire the system handles best (Generalisability), and whether a second run would return the same figure (Reproducibility).

\section{Limitations of LLMs in Truth-Seeking}\label{limitations-of-llms-in-truth-seeking}

Using a basic large language model in an evidence-retrieval setting has two key drawbacks. It can only recall information from its training data, and this information isn't necessarily true. LLMs are optimised to match the statistical distribution of their training corpus, so the gradient signal rewards predicting the most likely next token based on context, not factual truth. Truth only enters indirectly; the phrase ``the capital of France is'' is most likely followed by ``Paris'', as this is the most frequent resolution of the sentence, not because it `knows' that is the case. The training objective never targets truth directly; it is a byproduct of statistical patterns in text. Figure\,2.1 illustrates that, in the long tail where facts are rare in the training corpus, completing a statement may result in many continuations, not all of which are true (Kandpal et\,al.,\,2023). Backpropagation preserves the statistical regularity of a fact while discarding its provenance, so the model retains that a completion is likely but not where it came from or how reliable that source was.

\includegraphics[width=5.9in,height=2.26898in]{/Users/benjyb/Desktop/MscProject/.claude/worktrees/odmi-figures-reconciliation-244489/Dissertation/latex/media/media/image3.png}

\ccnote{Figure 2.1: Why the objective produces hallucination. }

A model built this way will therefore state things that are false, and will state them as confidently as it states things that are true. A pertinent question is why the behaviour survives post-training, where reinforcement learning shapes the model towards answering in a desired manner rather than merely fluently. Most benchmarks are graded binarily and award no credit for abstention, so a model answering "I do not know" scores identically to one answering incorrectly, while a model that guesses has positive expected value (Kalai et al., 2025). Optimisation against such benchmarks actively encourages confident fabrication under uncertainty, since the chance of guessing correctly is non-zero and the chance of scoring by abstaining is not. This incentive is inherited by any system built on such a model.

\emph{Extrinsic} hallucination is the main risk for the purposes of agentic research, where the generated output cannot be verified from the source (Maynez et al., 2020; Ji et al., 2023). Previous works in retrieval-augmented generation and web retrieval systems have constrained models to reason solely from web-derived evidence, treating parametric recall as inadmissible (Lewis et al., 2020; Nakano et al., 2021; Menick et al., 2022; Gao et al., 2023). \includegraphics[width=7.08611in,height=3.24167in]{/Users/benjyb/Desktop/MscProject/.claude/worktrees/odmi-figures-reconciliation-244489/Dissertation/latex/media/media/image4.png}

\ccnote{Figure 2.2: Prior-steered retrieval. }

The constraint is not airtight\textbf{\ccnote{, as the loop traced in Figure 2.2 shows, since}} parametric knowledge cannot be switched off and models tend to resolve conflicts between retrieved context and prior belief in favour of the prior (Longpre et al., 2021). \ccnote{The no-retrieval baseline is reported in Table F.1}, to isolate the effect of prior steering. The failure mode that dominates in this setting, however, comes from the ODMI's sparse evidence. Much of what the questionnaire asks about exists only in obscure national portals, in inconsistent formats, and in some cases nowhere public at all, which is the long-tail regime in which a model\textquotesingle s recall is least reliable (Mallen et al., 2023). A search over that space often returns a document on the same subject as the question, but that does not directly settle it. The agent then has to decide whether such a document supports its claim, and it decides under the incentive set out above, where committing scores and abstaining does not. The dominant failure is therefore over-reading, where the system treats related evidence as proof of the claim. Every citation it offers is real, so requiring the agent to cite its sources does not catch it.

A ``yes'' is an existential claim, and one document that establishes the practice discharges it, so retrieval is the right instrument for reaching it. A ``no'' on a question asking whether a portal feature exists is settled by inspection, because the search space is bounded and the agent can look at the portal and see. A ``no'' about a team\textquotesingle s internal practice is not, because the practice leaves no public trace, and finding nothing cannot separate a practice that is absent from evidence the agent did not reach.\\
\strut \\
§1.3 set this out as the move from the closed world a country reports under to the open world the agent searches (Reiter, 1978). The consequence for an agent under the guessing incentive runs one way. Weak evidence can be misinterpreted as a ``yes'', while little can be misinterpreted as a ``no'', so the system commits on weak support in one direction and withholds a correct answer in the other. This highlights the need for abstention, allowing the system to decline answering when the retrieved evidence is insufficient to support the claim.

Meeting the Selectivity standard set out in §2.2 requires three things of the architecture. The system has to be able to interrogate its own outputs, to subject them to adversarial challenge, and to return nothing at all where the evidence does not reach. The remainder of this chapter takes each in turn, \claudenote{self-verification in §2.4, adversarial multi-agent debate in §2.5 and abstention in §2.6}.

\section{Methods towards Verification}\label{methods-towards-verification}

\ccnote{Figure 2.3 traces the lineage between the past works, where no single work arrives at the set of requirements the ODMI requires.}

\includegraphics[width=5.4in,height=3.69858in]{/Users/benjyb/Desktop/MscProject/.claude/worktrees/odmi-figures-reconciliation-244489/Dissertation/latex/media/media/image5.png}

\ccnote{Figure 2.3: What each prior method adds and misses}

SelfCheckGPT (Manakul et al., 2023) samples the same response several times and measures how far the responses agree, on the finding that hallucinated content diverges under resampling while well-attested content does not. Dhuliawala et al. (2024) take an adjacent route, drafting an answer, breaking it into atomic verification questions, and answering each without the draft in view so the original claim cannot steer the check. They note that retrieval-based systems beat their architecture on very rare facts, which \claudenote{ties back to the long-tail argument in §2.3}. Both works improve accuracy over answering directly, and both share one error. They detect inconsistency within prior knowledge and say nothing about whether that prior knowledge is correct. Neither consults an external source, so both fail on Attributability with no evidence trail beyond the model\textquotesingle s own parameters.

Intrinsic self-correction fails to improve reasoning, and sometimes degrades it \ccnote{(Huang et al., 2024).} This is because the bottleneck is not error \emph{correction}, but error \emph{detection} (Tyen et al., 2024). If a model cannot detect its own error, then \claudenote{asking} the same model to check for it will miss the same blind spot. Furthermore, intrinsic self-correction may bolster \claudenote{the model's self-confidence} \ccnote{(Simhi et al., 2025)} in its answers regardless of whether the claim \ccnote{is true, degrading Selectivity}. This points towards needing an external knowledge base to fact-check claims before submission.

SAFE (Wei et al., 2024) decomposes a response into individual atomic claims, running an iterative search-and-reason loop over each claim and issuing its own follow-up queries before rating it ``supported'' or ``not supported''. This method passes Attributability, as it consults external evidence and \ccnote{ties each claim to a cited source. It fails} Selectivity because SAFE has no confidence output and no way to return nothing. The authors state that SAFE "does not assess the global accuracy of a fact" (p.28). A claim the search finds no evidence for is rated not supported, \ccnote{the same as a claim the search finds contradicting evidence for.} Absence of evidence becomes evidence of absence, which is the closed-world assumption §1.3 rules out.

The deeper limitation is the direction of the search. SAFE generates its queries from the claim in order to verify it, so the evidence it gathers is selected for proximity to the claim before anything is judged, and no step goes looking for material that would contradict it. The authors concede the consequence, that SAFE "may find supporting information for a factually inaccurate fact" (p.28). Nothing at the rating step corrects for this, because the model must return supported or not supported with no way to say it cannot tell. \claudenote{That is the binary grading §2.3 described}, where a model does better by guessing than by acknowledging uncertainty (Kalai et al., 2025). The ODMI setting may exacerbate this issue, because a question concerning obscure national data can be surrounded by ample publicly available adjacent evidence, yet fail to establish the data directly, and a verifier seeking only support may interpret the proximity as confirmation.

\ccnote{This motivates the question this work puts to test, whether opposition} is what a setting with little to no evidence requires. If one agent is tasked with defending a claim, and one with refuting it, the model's disposition to confidently assert is set against itself instead of being left to run in the same direction, meaning only claims backed by substantial \claudenote{evidence should survive the process}. This results in three relations. The claim survives refutation, the refutation defeats the claim, or both positions are asserted equally with opposing evidence. \ccnote{It is} this third relation \ccnote{that a single corroborative check cannot resolve and is most significant} in an environment where sparse and adjacent evidence is expected.\ccnote{ Whether the opposition itself does the work, or the extra retrieval that arguing forces, is what }\claudenote{§4.7}\ccnote{ separates.}

\section{Multi-Agent Debate for Verification}\label{multi-agent-debate-for-verification}

Multi-agent adversarial verification sits within the multi-agent debate family and is a relatively new line of research. Du et al. (2023) proposed a collaborative multi-agent architecture and found that cooperatively proposing and critiquing arguments meant the agents' collective intelligence could self-correct hallucinations. Liang et al. (2023)\claudenote{ are among the first to shift away from} collaborative debate, towards using adversarial tension to expose weak reasoning. \claudenote{§2.4 showed that a model cannot reliably check itself}, since it carries the same blind spots that produced the error (Huang et al., 2024; Tyen et al., 2024). A second agent set to agree inherits the problem and adds little in a data-scarce setting. Cohen et al. (2023) found that opposition surfaces errors a corroborative check misses, though their verifier interrogated the claim alone and admitted no external evidence. The hypothesis here is that adversarial pressure decorrelates the two models\textquotesingle{} errors, and that pairing that pressure with retrieval closes the gap Cohen left open.

If model A (the Researcher) is tasked with proving a claim, it is optimised to assert over abstaining (Kalai et al., 2025) and may retrieve some tangential evidence that proves an adjacent claim because no direct evidence exists on the web. Model B (the Verifier) is tasked with finding a flaw in model A's argument, and this weak reasoning should be surfaced. Note this decorrelation is only partial since A and B may share weights and priors, relevant here as this work uses the same Claude model for all agents. A and B argue their case over sequential rounds, gathering new evidence for each round, and if no agreement is reached this transcript is given to model C (the Adjudicator), which is comparatively reliable at correcting logical inconsistencies once their location is supplied (Tyen et al., 2024), as each agent has pointed out each other's weak reasoning.

Chowdhury et al. (2026) are the closest architectural precedent. Their PROClaim framework creates a judicial process for multi-agent debate, with evidence admission protocols and multi-stage deliberation, and it improves on standard multi-agent debate by 10.0 points. Each of their three judges runs a different underlying model, on the argument that agents drawn from the same weights share representations and cannot be relied on to disagree. Their own ablation puts the largest single share of the gain, 7.5 points, on progressive evidence retrieval. The argumentation structure accounts for less. Removing progressive retrieval raises agreement between their judges while accuracy falls. They read that rise as convergence without correctness, an ``epistemic bubble'', which is the failure a corroborative check produces on thin evidence. A single-call model with retrieval still beat the full courtroom on raw accuracy, 85.8\% against 81.7\%, however the same single model argued opposing conclusions convincingly from one evidence pool. An answer that cannot be audited is worth less than one that can, whatever it scores, and Chowdhury et al. rest their contribution on traceability and calibration rather than accuracy. They report an expected calibration error of 0.034, which addresses the second condition of Selectivity directly. It comes from a cross-validated grid search over one parameter, the weight on judge consensus in their confidence score, and their accuracy is unchanged at either setting, so the calibration rescales confidence without moving a verdict.

Smit et al. (2024) found that debate protocols do not reliably outperform cheaper strategies such as self-consistency and `multi-path ensembling'. Zhu et al. (2026) set out the sharper condition, showing that where agents are homogeneous and both are equally likely to be right, debate preserves expected correctness and cannot improve on majority vote. That applies here, since the Researcher and the Verifier are one model separated by a prompt. The result assumes a pool of evidence that stays fixed for the duration of the exchange, and each round here retrieves new material, which is the same mechanism Chowdhury et al. (2026) measured when progressive retrieval carried 7.5 of their 10.0 points. \ccnote{The claim under test is therefore narrower than ``debate is unanimously better''. It is that adversarial debate over retrieved evidence separates from a corroborative check in a sparse-web environment, whether on accuracy or on false positives. }\claudenote{§4.7}\ccnote{ settles it by ablating each agent in turn.}

\section{Allowing for Abstention}\label{allowing-for-abstention}

\claudenote{§2.5 left the system with three outcomes}. Debate can settle which of two cases has stronger evidence, but it cannot settle a case where the evidence is equally weak on both sides. Embedding a model in a debate \claudenote{does not remove the guessing incentive of §2.3}, and both Researcher and Verifier will guess if forced to produce an answer. Abstention removes this pressure. Where a question measures an internal practice, or little to no evidence exists online, abstention is what lets the system decline honestly.

One way to trigger it is to have the agent assign a confidence score to its argument and let a threshold decide when to abstain. It is much simpler to ask the agent to rate how well its retrieved evidence supports the claim, than it is to get it to refuse to answer outright due to insufficient evidence. Generating a confidence score demands the agent weigh up the relation between two pieces of text, whereas expecting it to self-regulate goes against its engrained guessing incentive. The decision to withhold then belongs to the threshold, and the agent is never asked to refuse.

When a model commits with confidence 0.7 it should be correct roughly 70\% of the time, but the score is self-declared, and nothing makes it track truth. Calibration is therefore assessed after the run, from the confidences the system logged. \ccnote{§3.9} defines the measure and \claudenote{§4.2} reports it. Feng et al. (2024) show calibrated abstention is achievable in multi-agent collaboration, reporting an abstain expected calibration error for their collaborative approaches. They also find models abstain less when predicting future election outcomes for African and Asian countries, which they read as a fairness concern for abstention mechanisms. Theirs is the only work in this chapter to address Selectivity and Subgroup Equity together.

Abstention has a hidden cost, as while the system declining is what stops it guessing, a declined question is still a question the assessment has no answer for. This means total accuracy reduces as abstention increases, and a system maximising accuracy just on answers it commits is incentivised to answer extremely cautiously. And thus, the choice of threshold is a policy decision, balancing coverage and accuracy. Selectivity from §2.2 asks both whether the system can decline and whether it was right to do so. The design supplies the first. Calibration is what would make the stated confidence non-trivial, and \claudenote{§4.2} reports whether it did.

\section{Multilingual Evidence Retrieval}\label{multilingual-evidence-retrieval}

The ODMI covers thirty-six states and more than twenty languages, and any variance in how well the system assesses them lands directly on Subgroup Equity. When a state scores poorly, three explanations sit behind that score, since the model may have misunderstood the evidence, the evidence may never have been published, or the country may simply be less open than its peers. Only the first is a property of the assessor rather than the country, which makes comprehension the one artefact that can and should be isolated.

Language capability is not uniform, because an LLM acquires each language as a by-product of next-token prediction and therefore in proportion to how often that language appeared in its corpus. Thellmann et al. (2024) evaluated forty LLMs on translated benchmarks across twenty-one European languages, and for LLaMA-3.1-70B-Instruct average accuracy falls from 0.79 in English to 0.66 or below for Estonian, Latvian and Lithuanian, declining with each language\textquotesingle s share of CommonCrawl's web archives. Anthropic\textquotesingle s own card shows the same trend inside the Claude family, with Sonnet scoring 69.0\% on the translated MMLU against 79.0\% on the English original (Anthropic, 2024).

Isolating comprehension does not clear Subgroup Equity on its own. Countries differ in what they publish, how they publish it, and how easy their national estates are to browse, so lower maturity may itself confound the ability for the system to assess their score. Separating them is what \claudenote{§4.4} attempts to do.

\section{Research Gap}\label{research-gap}

Two lines of work run through this chapter. One gives a model the means to retrieve its own evidence and act on it, the other checks whether an answer is sound. The two works that come closest to joining them are set out below, before \ccnote{Table 2.2} scores every work in the chapter against the six criteria.

Yao et al. (2023) interleave reasoning and action in their ReAct framework, alternating a thought, an action, and an observation of what that action returned, so that reasoning decides what to retrieve and what comes back shapes the next thought. The alternation leaves a legible record of each thought, the action it led to and what that action found, which is Attributability in principle. However, it does not provide verification. The loop ends once the agent judges it has an answer, so nothing in the architecture ensures the output is correct.

The closest analogue to this paper's work is Fumega and Gao (2026), who test commercial deep-research agents against the Global Data Barometer questionnaire, an assessment similar to the ODMI. A deep-research agent takes a question, runs its own searches, reads what it finds and returns a report with a list of sources. In their work, an agent finds the evidence itself and then answers from it, and they report an average match of 58\% against the human answers over four indicators. That runs from 81\% on Data Protection, whose evidence sits in a single legal document, to 31\% on Land Tenure, whose evidence is scattered across separate land management systems. They also identified cases in which the agent's answer was more accurate than the expert's, and whether a similar premise holds in the ODMI's case is taken up in \claudenote{§5.3}.

Their work highlights the necessity of both conditions of Attributability, as the agents are asked for citations and supply them, per sub-question with a working link, but nothing checks that the citation supports the claim. They cite such failures, e.g., an ``API Access'' button that did not function, as evidence without testing it. They have no independent computable measure to anchor Convergent Validity either, of the kind the nine computable ODMI questions provide. Subgroup Equity is raised and left untested, and there is no abstention channel or calibrated confidence.

Appendix E details how these works score against the full criteria.

\section{}\label{section-6}

{\def\LTcaptype{none} % do not increment counter
\begin{longtable}[]{@{}
  >{\raggedright\arraybackslash}p{(\linewidth - 12\tabcolsep) * \real{0.1639}}
  >{\raggedright\arraybackslash}p{(\linewidth - 12\tabcolsep) * \real{0.1382}}
  >{\raggedright\arraybackslash}p{(\linewidth - 12\tabcolsep) * \real{0.1394}}
  >{\raggedright\arraybackslash}p{(\linewidth - 12\tabcolsep) * \real{0.1393}}
  >{\raggedright\arraybackslash}p{(\linewidth - 12\tabcolsep) * \real{0.1415}}
  >{\raggedright\arraybackslash}p{(\linewidth - 12\tabcolsep) * \real{0.1395}}
  >{\raggedright\arraybackslash}p{(\linewidth - 12\tabcolsep) * \real{0.1382}}@{}}
\toprule\noalign{}
\begin{minipage}[b]{\linewidth}\raggedright
\textbf{Work}
\end{minipage} & \begin{minipage}[b]{\linewidth}\raggedright
\textbf{Conv. Validity}
\end{minipage} & \begin{minipage}[b]{\linewidth}\raggedright
\textbf{Generalis.}
\end{minipage} & \begin{minipage}[b]{\linewidth}\raggedright
\textbf{Subgroup Equity}
\end{minipage} & \begin{minipage}[b]{\linewidth}\raggedright
\textbf{Attributability}
\end{minipage} & \begin{minipage}[b]{\linewidth}\raggedright
\textbf{Selectivity}
\end{minipage} & \begin{minipage}[b]{\linewidth}\raggedright
\textbf{Reprod.}
\end{minipage} \\
\midrule\noalign{}
\endhead
\bottomrule\noalign{}
\endlastfoot
ReAct (Yao et al., 2023) & Partial & Partial & ✗ & Partial & ✗ & Partial \\
SelfCheckGPT (Manakul et al., 2023) & Partial & ✗ & ✗ & ✗ & Partial & Partial \\
Chain-of-Verification (Dhuliawala et al., 2024) & Partial & Partial & ✗ & ✗ & ✗ & Partial \\
LM vs LM (Cohen et al., 2023) & Partial & Partial & ✗ & ✗ & ✗ & Partial \\
Du et al. (2023) & Partial & Partial & ✗ & ✗ & ✗ & Partial \\
Liang et al. (2023) & Partial & ✗ & ✗ & ✗ & ✗ & Partial \\
Smit et al. (2024) & Partial & Partial & ✗ & ✗ & ✗ & Partial \\
Zhu et al. (2026) & Partial & Partial & ✗ & ✗ & Partial & Partial \\
SAFE (Wei et al., 2024) & ✓ & Partial & ✗ & ✓ & Partial & Partial \\
Feng et al. (2024) & Partial & Partial & ✓ & Partial & ✓ & Partial \\
Chowdhury et al. (2026) & ✓ & ✓ & ✗ & ✓ & ✓ & ✓ \\
Fumega and Gao (2026) & ✓ & ✓ & ✗ & Partial & ✗ & Partial \\
This work & ✓ & ✓ & ✓ & ✓ & ✓ & ✓ \\
\end{longtable}
}

\ccnote{Table 2.2}: Prior work scored against the six criteria of §2.2. A tick means the work addresses and reports the criterion, not that it succeeds. Partial means it addresses one of the criterion\textquotesingle s two conditions, or reports it without measuring it. The full justification for every cell, with the passage each verdict is read from, is in Appendix E.

\ccnote{Table 2.2} shows the shape of the gap. Attributability and Selectivity are met together only by Chowdhury et al. (2026). The only system operating in a real assessment setting, Fumega and Gao (2026), satisfies neither of the criteria fully. Although commercial research agents leave a per‑claim trail, they lack a method to verify that the cited source supports the claim, and they have no abstention channel. The verification machinery and the open-world setting have never been brought together.

Two columns of the table are close to empty. Only Chowdhury et al. (2026) test whether a generative assessor returns the same answers on a second run, reporting performance across three independent runs and conceding that run-level variance remains notable. Every other work states its decoding settings without ever repeating a run. Subgroup Equity has one entry. Feng et al. (2024) split their election set by region and found abstention rates lower for Africa and Asia, which they read as a fairness problem in abstention itself. That is the only prior measurement of whether such a system treats its subjects evenly, and \claudenote{§4.4} analyses how abstention varies within the subgroups.

This paper takes adversarial verification into that setting. A Verifier hunts for counterevidence, so a committed answer has had to survive an active attempt to refute it before an Adjudicator settles the case. Abstention is a first-class output, calibrated on confidence scores, so that thin evidence produces a withheld answer. The design targets the two criteria corroborative systems miss, namely Attributability, where every committed answer traces to a passage that entails it, and Selectivity, where stated confidence tracks truth and the system can decline. The remaining four criteria from §2.2 test whether the system succeeds, with Convergent Validity tested against the human key and against an independent recompute, Reproducibility upholding strict receipt taking and tracking stability across multiple runs, and Generalisability and Subgroup Equity looking across a decomposition by dimension and country.

